\chapter{Alloys and Reactive Bonds}\label{ch:alloys-bonds}

\index{alloys|(}
\index{reactive bonds|(}

So far, the phase system has been all-or-nothing: a value is either entirely fluid or entirely crystal.
Freeze an array, and every element freezes.
Freeze a struct, and every field hardens.
But real-world data is rarely so uniform.

Consider a user session object.
The user's ID and email should be immutable once established---nobody should accidentally overwrite them.
But the session's last-active timestamp needs to update on every request, and the permissions cache might need to refresh periodically.
You need a struct where \emph{some fields are crystal and others are fluid}---a material with mixed properties.

In materials science, such a material is called an \emph{alloy}: a combination of elements with different characteristics fused into a single structure.
Lattice borrows this metaphor.
An alloy struct has per-field phase declarations, so each field can have its own mutability.
And once you have mixed-phase data, Lattice goes further with \emph{reactive bonds}---automatic relationships between variables that propagate phase changes---and \emph{seeds}---validation contracts that guard the crystallization process.

\section{Alloy Structs: Per-Field Phase Declarations}\label{sec:alloy-structs}

\index{alloys!struct declarations}

An alloy struct declares the phase of each field individually using \lstinline|flux| and \lstinline|fix| annotations in the struct definition:

\begin{lstlisting}[language=Lattice, caption={Defining an alloy struct}]
struct Session {
    fix user_id: String,
    fix email: String,
    flux last_active: String,
    flux permissions: Array
}
\end{lstlisting}

When you create an instance of this struct, the runtime applies per-field phase enforcement.
The \lstinline|fix| fields are automatically frozen, while the \lstinline|flux| fields remain fluid:

\begin{lstlisting}[language=Lattice, caption={Creating an alloy struct instance}]
flux session = Session {
    user_id: "usr_42",
    email: "alice@example.com",
    last_active: "2024-01-15T10:30:00Z",
    permissions: ["read", "write"]
}

// Crystal fields are locked
print(session.user_id)     // "usr_42"
// session.user_id = "usr_99"  // Error: cannot assign to frozen field 'user_id'

// Fluid fields can change
session.last_active = "2024-01-15T10:31:00Z"
session.permissions.push("admin")
print(session.last_active)  // "2024-01-15T10:31:00Z"
\end{lstlisting}

\begin{defnbox}{Alloy}
\index{alloys!definition}
An \emph{alloy} is a struct (or map) with per-field phase declarations.
Individual fields carry their own phase tag, independent of the overall struct's phase.
Crystal fields reject mutation; fluid fields allow it.
The struct itself can be fluid (allowing field reassignment for fluid fields) while its crystal fields remain permanently locked.
\end{defnbox}

\subsection{How Alloys Work Under the Hood}

The per-field phase system is implemented through the \lstinline|field_phases| array in the \lstinline|LatValue| struct (see \filename{include/value.h}).
When a struct has per-field phases, \lstinline|field_phases| is a dynamically allocated array of \lstinline|PhaseTag| values---one for each field.
If \lstinline|field_phases| is \lstinline|NULL|, all fields inherit the struct's overall phase.

During struct construction in the VM (\filename{src/stackvm.c}), the runtime checks for a per-field phase declaration stored in the global environment under the key \lstinline|__struct_phases_<StructName>|.
This is an array of integer phase codes:

\begin{itemize}
  \item \lstinline|0| means \lstinline|PHASE_FLUID| --- the field stays mutable.
  \item \lstinline|1| means \lstinline|PHASE_CRYSTAL| --- the field is automatically frozen at construction.
  \item Any other value means the field inherits the struct's overall phase.
\end{itemize}

When field assignment is attempted (\lstinline|OP_SET_FIELD|), the VM checks the \lstinline|field_phases| array:

\begin{lstlisting}[language=Lattice, caption={Field-level phase enforcement}]
// Internally, the VM checks:
// if (obj.as.strct.field_phases[i] == VTAG_CRYSTAL)
//     ERROR: "cannot assign to frozen field 'user_id'"
\end{lstlisting}

This enforcement happens at the deepest level of the runtime---there is no way to bypass it through clever coding.
The field's phase tag is a property of the value itself, checked on every write.

\subsection{Partial Freeze with \texttt{except}}

\index{freeze@\texttt{freeze()}!partial freeze}

You can also create alloy-like behavior dynamically by freezing a struct with exceptions:

\begin{lstlisting}[language=Lattice, caption={Freezing a struct with exceptions}]
struct Config {
    host: String,
    port: Int,
    debug: Bool,
    log_level: String
}

flux config = Config {
    host: "localhost",
    port: 8080,
    debug: true,
    log_level: "info"
}

// Freeze everything except debug and log_level
freeze(config, except: ["debug", "log_level"])

// config.host = "evil.com"   // Error: cannot assign to frozen field 'host'
// config.port = 9999         // Error: cannot assign to frozen field 'port'
config.debug = false           // OK: debug was exempted
config.log_level = "debug"     // OK: log_level was exempted
\end{lstlisting}

The \lstinline|except| clause is compiled to the \lstinline|OP_FREEZE_EXCEPT| opcode.
Under the hood (in \filename{src/stackvm.c}), this opcode:

\begin{enumerate}
  \item Allocates the \lstinline|field_phases| array if it does not already exist.
  \item Iterates through each field, freezing it unless it appears in the exception list.
  \item Sets the exempted fields' phases to \lstinline|VTAG_FLUID|.
\end{enumerate}

The same mechanism works for maps---each key can have its own phase via the \lstinline|key_phases| map stored alongside the map's data.

\section{Designing Data with Mixed Mutability}\label{sec:designing-mixed-mutability}

\index{alloys!design patterns}

Alloy structs are a powerful design tool.
Here are patterns that emerge naturally from per-field phases.

\subsection{Identity + State Pattern}

Separate immutable identity from mutable state:

\begin{lstlisting}[language=Lattice, caption={Identity + State pattern}]
struct Player {
    fix id: Int,
    fix name: String,
    flux health: Int,
    flux position_x: Float,
    flux position_y: Float,
    flux inventory: Array
}

flux hero = Player {
    id: 1,
    name: "Aria",
    health: 100,
    position_x: 0.0,
    position_y: 0.0,
    inventory: ["sword", "shield"]
}

// Identity is locked forever
print(hero.name)  // "Aria"

// State changes freely
hero.health = hero.health - 15
hero.position_x = 3.5
hero.inventory.push("potion")
\end{lstlisting}

\subsection{Configuration + Runtime Pattern}

Lock configuration values while allowing runtime counters and caches to change:

\begin{lstlisting}[language=Lattice, caption={Configuration + Runtime pattern}]
struct Server {
    fix host: String,
    fix port: Int,
    fix max_connections: Int,
    flux active_connections: Int,
    flux request_count: Int
}

flux server = Server {
    host: "0.0.0.0",
    port: 443,
    max_connections: 10000,
    active_connections: 0,
    request_count: 0
}

// Configuration is safe from accidental changes
// Runtime counters update freely
server.active_connections = server.active_connections + 1
server.request_count = server.request_count + 1
\end{lstlisting}

\subsection{Append-Only Pattern}

Use a crystal array field for historical records (by replacing it with a new frozen array), while keeping a mutable accumulator:

\begin{lstlisting}[language=Lattice, caption={Append-only log pattern}]
struct AuditLog {
    fix entries: Array,
    flux pending: Array
}

fn add_entry(log: any, entry: any) {
    log.pending.push(entry)
    return log
}

fn flush_log(log: any) {
    flux all = thaw(log.entries)
    for item in log.pending {
        all.push(item)
    }
    // Return new log with frozen entries and empty pending
    return AuditLog {
        entries: freeze(all),
        pending: []
    }
}
\end{lstlisting}

\begin{tipbox}{Alloy Design Principle}
When designing an alloy struct, ask: ``Which fields define \emph{what this thing is} (identity), and which fields describe \emph{how it is right now} (state)?''
Identity fields should be \lstinline|fix|.
State fields should be \lstinline|flux|.
\end{tipbox}

\section{Reactive Bonds: \texttt{bond()}, \texttt{unbond()}, \texttt{react()}, \texttt{unreact()}}\label{sec:reactive-bonds}

\index{bond@\texttt{bond()}|(}
\index{react@\texttt{react()}|(}

Lattice provides a reactive system that lets variables respond automatically when the phase of another variable changes.
The core primitives are:

\begin{itemize}
  \item \lstinline|react(variable, callback)| --- register a callback that fires when a variable's phase changes.
  \item \lstinline|unreact(variable)| --- remove all reaction callbacks for a variable.
  \item \lstinline|bond(target, dependency, strategy)| --- create a dependency relationship between variables.
  \item \lstinline|unbond(target, dependency)| --- remove a specific dependency.
\end{itemize}

\subsection{Reactions: Responding to Phase Changes}

\lstinline|react()| registers a callback closure that is invoked whenever a variable undergoes a phase transition (freeze, thaw, or sublimate):

\begin{lstlisting}[language=Lattice, caption={Setting up a phase reaction}]
flux temperature = 72.5

react(temperature, |phase: any| {
    print("Temperature phase changed to: " + phase)
})

freeze(temperature)
// Output: Temperature phase changed to: crystal

thaw(temperature)
// Output: Temperature phase changed to: fluid
\end{lstlisting}

The callback receives a string describing the new phase: \lstinline|"crystal"|, \lstinline|"fluid"|, or \lstinline|"sublimated"|.
Multiple callbacks can be registered for the same variable---they fire in the order they were added.

Under the hood, reactions are stored in the \lstinline|LatRuntime| struct (see \filename{include/runtime.h}).
Each \lstinline|RTReaction| entry associates a variable name with an array of callback closures.
When \lstinline|OP_FREEZE_VAR| or \lstinline|OP_THAW_VAR| executes in the VM, it calls \lstinline|stackvm_fire_reactions()| to invoke all registered callbacks.

\begin{lstlisting}[language=Lattice, caption={Multiple reactions on one variable}]
flux sensor = 98.6

react(sensor, |phase: any| {
    print("Logger: sensor is now " + phase)
})

react(sensor, |phase: any| {
    if phase == "crystal" {
        print("Alert: sensor has been frozen!")
    }
})

freeze(sensor)
// Output: Logger: sensor is now crystal
// Output: Alert: sensor has been frozen!
\end{lstlisting}

\subsection{Removing Reactions}

\lstinline|unreact()| removes all reaction callbacks registered on a variable:

\begin{lstlisting}[language=Lattice, caption={Removing all reactions}]
flux value = 42
react(value, |phase: any| {
    print("phase changed: " + phase)
})

unreact(value)
freeze(value)  // No output---the callback has been removed
\end{lstlisting}

\subsection{Bonds: Phase Dependency Relationships}

\index{bond@\texttt{bond()}!dependency}

While reactions let you \emph{observe} phase changes, bonds let you \emph{propagate} them.
\lstinline|bond()| creates a dependency relationship: when a dependency variable freezes, the target variable also freezes automatically.

\begin{lstlisting}[language=Lattice, caption={Creating a bond between variables}]
flux primary = [1, 2, 3]
flux replica = [1, 2, 3]

bond(replica, "primary", "mirror")

// When primary freezes, replica automatically freezes too
freeze(primary)
print(phase_of(primary))  // "crystal"
print(phase_of(replica))  // "crystal" (frozen via cascade)
\end{lstlisting}

The third argument to \lstinline|bond()| is the \emph{strategy} string, which determines how the freeze cascades.
The \lstinline|"mirror"| strategy (the default) means the target mirrors the dependency's phase transitions.

\begin{notebox}{Bond Requirements}
Bonds have two requirements:
\begin{enumerate}
  \item The dependency variable must exist at the time of bonding.
  \item The target variable must not already be frozen---you cannot bond a crystal variable (it is already locked).
\end{enumerate}
Violating either rule produces a runtime error: ``\texttt{cannot bond undefined variable}'' or ``\texttt{cannot bond already-frozen variable}.''
\end{notebox}

\subsection{Removing Bonds}

\lstinline|unbond()| removes a specific dependency from a target:

\begin{lstlisting}[language=Lattice, caption={Removing a bond}]
flux source = 100
flux follower = 100

bond(follower, "source", "mirror")
unbond(follower, "source")

freeze(source)
print(phase_of(follower))  // "fluid" (no longer bonded)
\end{lstlisting}

If removing a dependency leaves a bond entry with no dependencies, the entire bond entry is cleaned up.

\index{bond@\texttt{bond()}|)}
\index{react@\texttt{react()}|)}

\section{Bond Strategies: Mirror, Inverse, Gate}\label{sec:bond-strategies}

\index{bond strategies|(}

The strategy parameter in \lstinline|bond()| controls how phase transitions propagate from dependency to target.

\subsection{Mirror Strategy}

\index{bond strategies!mirror}

The \lstinline|"mirror"| strategy is the default.
When the dependency freezes, the target freezes.
When the dependency thaws, the target thaws.
The target mirrors whatever the dependency does:

\begin{lstlisting}[language=Lattice, caption={Mirror bond strategy}]
flux leader = "hello"
flux follower = "world"

bond(follower, "leader", "mirror")

freeze(leader)
print(phase_of(follower))  // "crystal" (mirrored freeze)
\end{lstlisting}

The mirror strategy is implemented in \lstinline|rt_freeze_cascade()| in \filename{src/runtime.c}.
When a variable is frozen, the cascade function walks all bonds, finds targets whose dependencies include the frozen variable, and applies the appropriate phase transition.

\subsection{Designing Custom Strategies}

Bond strategies are stored as string identifiers in the \lstinline|RTBond| struct's \lstinline|dep_strategies| array.
The runtime matches on these strings during cascade operations.
While the core runtime ships with the mirror strategy, the design is extensible---additional strategies can be recognized by the cascade logic.

Here are patterns you can build with the existing bond primitives:

\subsubsection{Inverse Pattern}

Using reactions, you can implement an inverse relationship where freezing one variable thaws another:

\begin{lstlisting}[language=Lattice, caption={Inverse relationship via reactions}]
flux active_config = Map::new()
active_config.set("theme", "dark")

flux backup_config = Map::new()
backup_config.set("theme", "dark")

react(active_config, |phase: any| {
    if phase == "crystal" {
        // When active freezes, make a thawed backup
        print("Active config frozen, backup is now active")
    }
})

freeze(active_config)
// Output: Active config frozen, backup is now active
\end{lstlisting}

\subsubsection{Gate Pattern}

A gate pattern uses a boolean sentinel to control whether a freeze cascades:

\begin{lstlisting}[language=Lattice, caption={Gate pattern with reactions}]
flux gate_open = true
flux protected_data = [1, 2, 3]

react(protected_data, |phase: any| {
    if phase == "crystal" {
        print("Data has been sealed")
    }
})

fn conditional_freeze(data: any, gate: any) {
    if gate {
        return freeze(data)
    }
    return data
}

// With gate open, freeze proceeds
fix sealed = conditional_freeze(protected_data, gate_open)
// Output: Data has been sealed
\end{lstlisting}

\index{bond strategies|)}

\section{Pressure: \texttt{pressurize()}, \texttt{depressurize()}, \texttt{pressure\_of()}}\label{sec:pressure}

\index{pressure|(}
\index{pressurize@\texttt{pressurize()}|(}

Pressure is a constraint system for \emph{structural mutations} on collections.
While the phase system controls whether a value can be mutated at all, pressure controls \emph{how} it can be mutated---specifically, whether a collection can grow, shrink, or both.

\begin{lstlisting}[language=Lattice, caption={Applying pressure constraints}]
flux log_entries = ["boot", "init"]

// Apply pressure: no shrinking allowed
pressurize("log_entries", "no_shrink")

log_entries.push("request received")  // OK: growing is allowed
print(log_entries)  // ["boot", "init", "request received"]

// log_entries.pop()  // Error: pressure constraint violated (no_shrink)
\end{lstlisting}

\subsection{Pressure Modes}

The \lstinline|pressurize()| function accepts a variable name (as a string) and a mode string.
Four modes are available:

\begin{table}[h]
\centering
\begin{tabular}{lp{8cm}}
\toprule
\textbf{Mode} & \textbf{Effect} \\
\midrule
\lstinline|"no_grow"| & Blocks \lstinline|push()|, \lstinline|insert()|, and other size-increasing operations. \\
\lstinline|"no_shrink"| & Blocks \lstinline|pop()|, \lstinline|remove_at()|, and other size-reducing operations. \\
\lstinline|"no_resize"| & Blocks both growing and shrinking---the collection's size is locked. \\
\lstinline|"read_heavy"| & An optimization hint indicating the collection will be read much more often than written. \\
\bottomrule
\end{tabular}
\caption{Pressure modes and their effects.}\label{tab:pressure-modes}
\end{table}

\begin{lstlisting}[language=Lattice, caption={Different pressure modes}]
flux scores = [85, 92, 78, 95]

// Lock the size entirely
pressurize("scores", "no_resize")

// scores.push(100)      // Error: pressure constraint (no_resize)
// scores.pop()           // Error: pressure constraint (no_resize)
scores[0] = 90            // OK: element modification is still allowed
print(scores)  // [90, 92, 78, 95]
\end{lstlisting}

\begin{notebox}{Pressure vs.\ Phase}
Pressure and phase are orthogonal.
A fluid value under \lstinline|"no_resize"| pressure can still have its \emph{elements} modified---you just cannot add or remove elements.
A crystal value rejects \emph{all} modifications regardless of pressure.
Think of pressure as a structural constraint (shape of the container) and phase as a material constraint (mutability of the contents).
\end{notebox}

\subsection{Querying and Removing Pressure}

Use \lstinline|pressure_of()| to check the current pressure on a variable, and \lstinline|depressurize()| to remove it:

\begin{lstlisting}[language=Lattice, caption={Querying and removing pressure}]
flux items = [1, 2, 3]

pressurize("items", "no_grow")
print(pressure_of("items"))  // "no_grow"

depressurize("items")
print(pressure_of("items"))  // nil (no pressure)

items.push(4)  // Now allowed again
print(items)   // [1, 2, 3, 4]
\end{lstlisting}

\subsection{Implementation Details}

Pressure constraints are stored in the \lstinline|LatRuntime| struct as an array of \lstinline|RTPressure| entries (see \filename{include/runtime.h}).
Each entry maps a variable name to a mode string.
The pressure check is performed by array mutating methods in the VM (e.g., during \lstinline|push()|, \lstinline|pop()|, \lstinline|insert()|, and \lstinline|remove_at()|), which look up the variable name in the pressure table and reject operations that violate the constraint.

The native functions \lstinline|native_pressurize()|, \lstinline|native_depressurize()|, and \lstinline|native_pressure_of()| are defined in \filename{src/runtime.c}.
They validate the mode string, update the pressure table, and handle the swap-remove pattern for deregistration.

\index{pressurize@\texttt{pressurize()}|)}
\index{pressure|)}

\section{Seeds: \texttt{seed()}, \texttt{unseed()}}\label{sec:seeds}

\index{seed@\texttt{seed()}|(}

Seeds are \emph{validation contracts} that guard the crystallization process.
When you seed a variable, you attach a closure that must approve the value before it can be frozen.
If the seed contract rejects the value, the freeze fails with an error.

\begin{lstlisting}[language=Lattice, caption={Setting a seed contract}]
flux user_age = 25

seed(user_age, |val: any| {
    if val < 0 {
        return "age cannot be negative"
    }
    if val > 150 {
        return "age exceeds reasonable limit"
    }
    return nil  // nil means "approved"
})

// This freeze works: 25 passes the contract
fix sealed_age = freeze(user_age)
print(sealed_age)  // 25
\end{lstlisting}

\begin{lstlisting}[language=Lattice, caption={Seed contract rejecting a freeze}]
flux invalid_score = -10

seed(invalid_score, |val: any| {
    if val < 0 {
        return "score must be non-negative"
    }
    return nil
})

// This freeze fails: -10 does not pass the contract
// fix bad = freeze(invalid_score)
// Error: score must be non-negative
\end{lstlisting}

\begin{defnbox}{Seed}
\index{seed@\texttt{seed()}!definition}
A \emph{seed} is a validation contract (a closure) attached to a variable.
The seed is checked when the variable is frozen via \lstinline|freeze()| or \lstinline|grow()|.
If the closure returns a non-nil value (typically an error message string), the freeze is rejected.
If it returns \lstinline|nil|, the freeze proceeds.
\end{defnbox}

\subsection{Seeds and \texttt{grow()}}

\index{grow@\texttt{grow()}}

The \lstinline|grow()| function is a higher-level operation that validates all seeds and then freezes the variable.
Unlike a bare \lstinline|freeze()|, \lstinline|grow()| \emph{consumes} the seed contracts---after a successful \lstinline|grow()|, the seeds are removed:

\begin{lstlisting}[language=Lattice, caption={Using grow() with seeds}]
flux balance = 1000

seed(balance, |val: any| {
    if val < 0 { return "balance cannot be negative" }
    return nil
})

// grow() validates seeds, freezes, and removes the seed contracts
fix final_balance = grow("balance")
print(final_balance)  // 1000
// The seed contract is now consumed
\end{lstlisting}

The \lstinline|grow()| function is implemented as \lstinline|native_grow()| in \filename{src/runtime.c}.
It calls \lstinline|rt_validate_seeds()| with the \lstinline|consume| flag set to \lstinline|true|, which both validates the contracts and removes them from the runtime's seed table upon success.
After validation, it freezes the value, records the phase change in the tracking history, fires any cascade bonds, and triggers reactions.

\subsection{Removing Seeds Manually}

\lstinline|unseed()| removes the seed contract for a variable without triggering any validation:

\begin{lstlisting}[language=Lattice, caption={Removing a seed contract}]
flux data = "sensitive"

seed(data, |val: any| {
    if len(val) < 8 { return "too short" }
    return nil
})

// Changed our mind---remove the seed
unseed(data)

// Now freeze works without validation
fix sealed = freeze(data)
\end{lstlisting}

\subsection{Multiple Seeds}

You can attach multiple seed contracts to the same variable.
All of them must approve the value for the freeze to succeed:

\begin{lstlisting}[language=Lattice, caption={Multiple seed contracts}]
flux password = "hunter2"

seed(password, |val: any| {
    if len(val) < 8 { return "password must be at least 8 characters" }
    return nil
})

seed(password, |val: any| {
    if val == "password" { return "password cannot be 'password'" }
    return nil
})

// "hunter2" is only 7 characters---first seed rejects it
// fix sealed = freeze(password)
// Error: password must be at least 8 characters
\end{lstlisting}

\section{Building Reactive Data Flows}\label{sec:reactive-data-flows}

\index{reactive data flows}

By combining bonds, reactions, seeds, and pressure, you can build sophisticated reactive data flows where phase transitions cascade through your data model automatically.

\subsection{Example: Configuration Validation Pipeline}

Here is a complete example that uses seeds to validate configuration, bonds to cascade freezing from a primary config to dependent caches, and reactions to log all phase changes:

\begin{lstlisting}[language=Lattice, caption={A reactive configuration pipeline}]
// Set up the primary configuration
flux app_config = Map::new()
app_config.set("database_url", "postgres://localhost/mydb")
app_config.set("cache_ttl", "300")
app_config.set("max_retries", "3")

// Set up a dependent cache configuration
flux cache_config = Map::new()
cache_config.set("ttl", "300")
cache_config.set("backend", "redis")

// Bond: when app_config freezes, cache_config freezes too
bond(cache_config, "app_config", "mirror")

// Seed: validate configuration before freezing
seed(app_config, |val: any| {
    if val.get("database_url") == nil {
        return "database_url is required"
    }
    return nil
})

// Reaction: log phase changes
react(app_config, |phase: any| {
    print("app_config phase -> " + phase)
})

react(cache_config, |phase: any| {
    print("cache_config phase -> " + phase)
})

// Freeze the primary config
// 1. Seed validates (database_url exists) -> passes
// 2. app_config freezes
// 3. Reaction fires: "app_config phase -> crystal"
// 4. Bond cascade: cache_config freezes
// 5. Reaction fires: "cache_config phase -> crystal"
fix sealed_config = freeze(app_config)

print(phase_of(app_config))    // "crystal"
print(phase_of(cache_config))  // "crystal"
\end{lstlisting}

\subsection{Example: Sensor Data with Pressure Guards}

\begin{lstlisting}[language=Lattice, caption={Sensor data with pressure and seeds}]
flux readings = []

// Pressure: this buffer only grows (no accidental data loss)
pressurize("readings", "no_shrink")

// Seed: require at least 10 readings before allowing freeze
seed(readings, |val: any| {
    if len(val) < 10 {
        return "need at least 10 readings before sealing"
    }
    return nil
})

// Collect sensor data
for i in 0..15 {
    readings.push(20.0 + to_float(i) * 0.5)
}

// readings.pop()  // Error: pressure constraint (no_shrink)

// Now we have 15 readings, seed contract will pass
fix sealed_readings = freeze(readings)
print(sealed_readings.len())  // 15
\end{lstlisting}

\subsection{Example: Multi-Stage Initialization}

\begin{lstlisting}[language=Lattice, caption={Multi-stage initialization with bonds}]
flux database = Map::new()
flux http_server = Map::new()
flux app = Map::new()

// Set up cascading bonds
bond(http_server, "database", "mirror")
bond(app, "http_server", "mirror")

// Reactions for monitoring
react(database, |phase: any| {
    print("Database: " + phase)
})
react(http_server, |phase: any| {
    print("HTTP Server: " + phase)
})
react(app, |phase: any| {
    print("Application: " + phase)
})

// Configure each component
database.set("host", "localhost")
database.set("port", "5432")

http_server.set("port", "8080")
http_server.set("workers", "4")

app.set("name", "MyApp")
app.set("version", "1.0.0")

// Freeze the database layer
// This cascades: database -> http_server -> app
freeze(database)
// Output:
//   Database: crystal
//   HTTP Server: crystal
//   Application: crystal

print(phase_of(app))  // "crystal"
\end{lstlisting}

\begin{warnbox}{Cascade Ordering}
Bond cascades propagate synchronously during the freeze operation.
If you have a chain of bonds (A $\to$ B $\to$ C), freezing A will freeze B, which will then freeze C, all within the same operation.
Be careful with long chains---each step involves a full freeze (deep clone into an arena).
\end{warnbox}

\section{Exercises}\label{sec:ch13-exercises}

\begin{enumerate}
  \item \textbf{Alloy Design.}
    Design an alloy struct for a bank account with immutable fields for the account number and holder name, and mutable fields for the balance and transaction history.
    Write functions to deposit, withdraw, and display the account.
    Ensure the account number can never be changed.

  \item \textbf{Reactive Counter.}
    Create two variables, \lstinline|counter| and \lstinline|display|.
    Set up a reaction on \lstinline|counter| that prints the phase change, and a bond from \lstinline|display| to \lstinline|counter| using the mirror strategy.
    Freeze \lstinline|counter| and verify that \lstinline|display| freezes automatically.

  \item \textbf{Seed Validation.}
    Write a seed contract for an email address variable that validates the value contains an \lstinline|@| character and has at least 5 characters.
    Test the seed by trying to freeze valid and invalid email addresses.

  \item \textbf{Pressure Experiment.}
    Create an array under \lstinline|"no_grow"| pressure.
    Verify that \lstinline|push()| fails but element assignment works.
    Then depressurize and confirm \lstinline|push()| works again.
    Switch to \lstinline|"no_resize"| and verify both \lstinline|push()| and \lstinline|pop()| fail.

  \item \textbf{Complete Pipeline.}
    Build a complete reactive pipeline with three variables: \lstinline|raw_data|, \lstinline|validated_data|, and \lstinline|archived_data|.
    Use seeds to validate \lstinline|raw_data| before freezing, bonds to cascade the freeze to \lstinline|validated_data| and then \lstinline|archived_data|, and reactions to log every phase transition.
\end{enumerate}

\bigskip

\noindent\textbf{What's Next.}\quad
We have explored the full range of Lattice's phase system: from basic \lstinline|flux|/\lstinline|fix| declarations through alloy structs, reactive bonds, and validation seeds.
But all of this has operated in \emph{casual} mode, where the runtime catches phase violations dynamically.
In \cref{ch:strict-mode}, we explore strict mode and the static phase checker---a compile-time analysis that catches phase errors before your code ever runs.

\index{alloys|)}
\index{reactive bonds|)}
